\documentclass[10pt,a4paper]{amsart}
\usepackage[margin=19mm]{geometry}
\usepackage{amsmath,amssymb,mathtools}
\usepackage[scr=rsfs,cal=euler]{mathalfa}
\usepackage{xfrac}
\usepackage[unicode,hidelinks,bookmarks=false]{hyperref}
\newcommand{\ie}{i.e.\ }
\newcommand{\cf}{cf.\ }
\newcommand{\resp}{respectively\ }
\newcommand{\Subsection}{\subsection}
\providecommand{\nobreakdash}{\nobreak\hbox{-}}
\setlength{\emergencystretch}{2em}
\multlinegap0pt
\usepackage{bm}
\usepackage{bbm}
\usepackage{dsfont}
\usepackage{xspace}
\usepackage{cancel}
\usepackage{mathtools}
\usepackage{cleveref}
\usepackage{amsthm}
\makeatletter
\@ifundefined{assumption}{\newtheorem{assumption}{Assumption}}{}
\@ifundefined{lemma}{\newtheorem{lemma}{Lemma}}{}
\makeatother
\makeatletter
\newcommand{\RR}{\mathbb{R}}
\newcommand{\bD}{\bm{D}}
\newcommand{\bS}{\bm{S}}
\newcommand{\ba}{\bm{a}}
\newcommand{\bb}{\bm{b}}
\newcommand{\bc}{\bm{c}}
\newcommand{\be}{\bm{e}}
\newcommand{\bv}{\bm{v}}
\newcommand{\bw}{\bm{w}}
\newcommand{\bx}{\bm{x}}
\newcommand{\bz}{\bm{z}}
\newcommand{\calL}{\mathcal{L}}
\newcommand{\calM}{\mathcal{M}}
\expandafter\ifx\csname proofsketch\endcsname\relax
\newenvironment{proofsketch}
    {\noindent\textit{Proof sketch.}}
\fi
\newcommand{\prox}{\mathrm{prox}}
\makeatother
\title[MoSSP: A Momentum-Based Single-Loop Stochastic Penalty Metho]{MoSSP: A Momentum-Based Single-Loop Stochastic Penalty Method for Nonconvex Constrained DC-Regularized Optimization}
\author{Luxuan Li, Chunfeng Cui, Xiao Wang}
\date{Source: arXiv:2605.29635v1 (CC BY 4.0)}
\begin{document}
\maketitle

\noindent\emph{Source and licence.} MoSSP: A Momentum-Based Single-Loop Stochastic Penalty Method for Nonconvex Constrained DC-Regularized Optimization. Version arXiv:2605.29635v1 (CC BY 4.0), distributed under the Creative Commons Attribution 4.0 International licence, as recorded in the arXiv licence metadata for this exact version and shown on the abstract page https://arxiv.org/abs/2605.29635v1. Full rights notice: licenses/arxiv-2605.29635-CC-BY-4.0.txt.\par\smallskip

\noindent\textit{Editorial scope.} Counted unit: the Lemma whose source label is \texttt{Lemma: approximate-sufficient-descent} in the article (source0.tex lines 1390--1413, source label \texttt{Lemma: approximate-sufficient-descent}), delivered together with the article's own complete original proof of it (source0.tex lines 1415--1500, one \texttt{proof} environment of 86 lines). No mathematics is written, repaired or re-derived: statement and proof are the article's own text line for line. Required context retained uncounted: Ass: boundedness of c and f (lines 489--499); eq:bz\_update (lines 562--565); eq:bx\_update\_Polyak (lines 596--599); eq:bx\_update\_storm (lines 735--739); Eq: basic parameter conditions for both algorithms (lines 1382--1388). Every definition, display and label of those lines is retained unchanged. Omitted from this package and not redistributed: 1--488, 500--561, 566--595, 600--734, 740--1381, 1389--1389, 1414--1414, 1501--2243. Nothing inside a retained excerpt is omitted or altered: every retained body is delivered exactly as the article's lines are written, so no omission receipt of any kind accompanies this package. Citations: the retained excerpts carry no citation command at all, so no bibliography entry of the article is transcribed and no reference is redistributed. Reference adaptation: none was needed and none was made; every reference inside the delivered unit resolves inside it, so no unresolved reference or citation is produced. Printed slips kept literal: no printed word, symbol, hypothesis or number of the counted statement or of its proof was altered. Layout and class: the article's own class is \texttt{article} with packages including microtype, graphicx, subcaption, booktabs, hyperref, multirow, icml2026, amsmath, amssymb, mathtools, amsthm, bm, cleveref; the wrapper is the lane's pinned \texttt{amsart} editorial wrapper at 10pt on A4, which restates the theorem-like environments the retained text needs and adds the article's own macro definitions that the retained text needs (\texttt{\textbackslash RR}, \texttt{\textbackslash bD}, \texttt{\textbackslash bS}, \texttt{\textbackslash ba}, \texttt{\textbackslash bb}, \texttt{\textbackslash bc}, \texttt{\textbackslash be}, \texttt{\textbackslash bv}, \texttt{\textbackslash bw}, \texttt{\textbackslash bx}, \texttt{\textbackslash bz}, \texttt{\textbackslash calL}, \texttt{\textbackslash calM}, \texttt{\textbackslash prox}), with extra wrapper packages (bm, bbm, dsfont, xspace, cancel, mathtools, cleveref). The delivered numbering is the unit's own. Assets: no retained span contains a figure environment, a graphics-inclusion command, a PDF-page inclusion, a table or a TikZ picture; no image file is copied into this package, the package declares no asset dependency and no image file is redistributed with it. Licence: version arXiv:2605.29635v1 of this article is distributed under the Creative Commons Attribution 4.0 International licence, as recorded in the arXiv licence metadata for this exact version and shown on the abstract page https://arxiv.org/abs/2605.29635v1. A full rights notice is delivered at licenses/arxiv-2605.29635-CC-BY-4.0.txt. Characters: every retained excerpt is pure ASCII, so the delivered engine, XeLaTeX with its default Latin Modern text font, reports no missing character.
\begin{assumption}
    \label{Ass: boundedness of c and f} Function $f(\cdot)$ is
    $L_f$-smooth. Function \(\bc\) is \(L_c\)-smooth. There exist $C, G>0$ such that
    \begin{align}
         &\|\nabla f(\bx) \| \leq G,\,
         \sup_{\bv_h\in\partial h(\bx)}\|\bv_h\| \leq G,\,
         \sup_{\bv_g\in\partial g(\bx)}\|\bv_g\| \leq G,\\ \notag
         &\|\nabla \bc(\bx) \| \leq G,\quad
         \| \bc(\bx) \| \leq C,\quad \forall \bx\in\RR^n.\notag
    \end{align}
\end{assumption}

\begin{align} 
\label{eq:bz_update}
\bz^{k+1} = \bz^k - \beta \left( \prox_{\mu_k g}(\bz^k) - \bx^{k+1} \right),
\end{align}

\begin{align}
\bx^{k+1} = 
\prox_{\mu_k h}\bigl(\bz^k - \mu_k \bS^{k}\bigr), \label{eq:bx_update_Polyak}
\end{align}

\begin{align}
\label{eq:bx_update_storm}
    \bx^{k+1} 
    = \prox_{\mu_k h}\bigl(\bz^k - \mu_k \bD^{k}\bigr).
\end{align}

\begin{align}
    \label{Eq: basic parameter conditions for both algorithms}
    &\mu_{k}L_{\rho_k} \leq \frac{1}{4}, \quad
    \mu_{k+1} \leq \mu_{k}, \quad
    \rho_{k} \leq \rho_{k+1}, \quad
    0 < \beta \leq 1, \quad \forall\, k \geq 0.
\end{align}

\begin{lemma}
    \label{Lemma: approximate-sufficient-descent}
    Suppose that \Cref{Ass: boundedness of c and f}
    holds, and the parameters $\mu_{k}$, $\rho_{k}$ and $\beta$ satisfy \eqref{Eq: basic parameter conditions for both algorithms}.
    Then, for any $k \geq 0$, it holds that
    \begin{align}
        \label{eq:approximate-sufficient-descent}
        \calL_{\rho_{k+1}, \mu_{k+1}}(\bw^{k+1}) \leq&\, \calL_{\rho_k, \mu_k}(\bw^{k}) - \frac{(2\mu_{k})^{-1}- L_{\rho_k}}{2}\| \bw^{k+1}- \bw^{k}\|^{2} + \frac{\rho_{k+1}- \rho_{k}}{2}C^{2} + \Delta_{k+1} +  \mu_{k}\| \be^{k}\|^{2},
    \end{align}
where $\Delta_{k+1}
:= \frac{|\mu_{k}-\mu_{k+1}|}{2\mu_{k+1}^2}
(C^{2} + \|\bx^{k+1}-\bz^{k+1}\|^{2})$. Moreover, it holds that
\begin{align}
\label{eq:sequence-residual-bound}
\| \bw^{k+1}- \bw^{k}\|^{2}
\leq&\, \frac{4 \mu_{k}}{1 - 2 \mu_{k}L_{\rho_k}}
   \bigl(
      \calL_{\rho_k, \mu_k}(\bw^{k})
      - \calL_{\rho_{k+1}, \mu_{k+1}}(\bw^{k+1})
      + \Delta_{k+1} + \frac{\rho_{k+1}- \rho_{k}}{2}C^{2}
      + \mu_{k}\| \be^{k}\|^{2}
   \bigr).
\end{align}
\end{lemma}

\begin{proof}
    Since $\mu_{k}L_{\rho_k}\leq \frac{1}{4}$, it follows that
    $\frac{(2\mu_{k})^{-1}- L_{\rho_k}}{2}= \frac{1}{4 \mu_{k}}- \frac{L_{\rho_k}}{2}
    = \frac{1 - 2\mu_k L_{\rho_k}}{4\mu_k}
    \geq \frac{1}{8 \mu_{k}}> 0$.
    To prove \eqref{eq:approximate-sufficient-descent},
    we bound $\calL_{\rho_{k+1}, \mu_{k+1}}(\bw^{k+1}) - \calL_{\rho_k, \mu_k}(\bw^{k})$ by decomposing
    it as
    \begin{flalign}
        &\calL_{\rho_{k+1}, \mu_{k+1}}(\bw^{k+1}) - \calL_{\rho_k,\mu_k}(\bw^{k}) = [ \calL_{\rho_{k+1}, \mu_{k+1}}(\bw^{k+1}) - \calL_{\rho_k, \mu_k}(\bw^{k+1}) ]  + [ \calL_{\rho_k, \mu_{k}}(\bw^{k+1}) - \calL_{\rho_k, \mu_k}(\bw^{k}) ]. \label{eq:function-value-decomposition}
    \end{flalign}

For the first term in \eqref{eq:function-value-decomposition}, we have
\begin{align}
\calL_{\rho_{k+1}, \mu_{k+1}}(\bw^{k+1}) - \calL_{\rho_k, \mu_k}(\bw^{k+1})  \leq \frac{\rho_{k+1}- \rho_{k}}{2}C^{2} + \Delta_{k+1}, \label{eq:penalty-term-bound}
\end{align}
where we use $\|\bc(\bx^{k+1})\|\le C$ from \Cref{Ass: boundedness of c and f}.

    For the second term of \eqref{eq:function-value-decomposition}, we decompose it as
    \begin{flalign}
        &\calL_{\rho_k, \mu_k}(\bw^{k+1}) - \calL_{\rho_k, \mu_k}(\bw^{k}) \notag \\
        =&\, [ \calL_{\rho_k, \mu_k}(\bx^{k+1}, \bz^{k+1}) - \calL_{\rho_k, \mu_k}(\bx^{k+1}, \bz^{k}) ] + [ \calL_{\rho_k, \mu_k}(\bx^{k+1}, \bz^{k}) - \calL_{\rho_k, \mu_k}(\bx^{k}, \bz^{k}) ]. \label{eq:objective-function-decomposition}
    \end{flalign}
    The first part of \eqref{eq:objective-function-decomposition} is
\begin{align}
&\calL_{\rho_k,\mu_k}(\bx^{k+1},\bz^{k+1})
 - \calL_{\rho_k,\mu_k}(\bx^{k+1},\bz^{k}) \notag \\
=&\, \frac{1}{2\mu_k}(\|\bx^{k+1}-\bz^{k+1}\|^{2}
    - \|\bx^{k+1}-\bz^{k}\|^{2}) +  \calM_{\mu_k g}(\bz^{k}) - \calM_{\mu_k g}(\bz^{k+1}) \notag \\
\leq&\,\frac{1}{2\mu_k}(\|\bx^{k+1}-\bz^{k+1}\|^{2}
    - \|\bx^{k+1}-\bz^{k}\|^{2})  - \frac{1}{\mu_k}\langle \bz^{k}-\prox_{\mu_k g}(\bz^{k}),
     \bz^{k+1}-\bz^{k}\rangle. \label{eq:z-first-part}
\end{align}
Using the identity
\(
\|\ba\|^{2}-\|\bb\|^{2}
= 2\langle \ba-\bb,\ba\rangle - \|\ba-\bb\|^{2},
\)
with $\ba=\bx^{k+1}-\bz^{k+1}$ and $\bb=\bx^{k+1}-\bz^{k}$, we obtain
\begin{align*}
\|\bx^{k+1}-\bz^{k+1}\|^{2}-\|\bx^{k+1}-\bz^{k}\|^{2}
= 2\langle \bz^{k}-\bz^{k+1},\bx^{k+1}-\bz^{k+1}\rangle
  - \|\bz^{k+1}-\bz^{k}\|^{2}.
\end{align*}
Substituting this into \eqref{eq:z-first-part} and combining the optimality condition of the subproblem \eqref{eq:bz_update},
\begin{align}
\label{Eq: the optimality of z}
    \prox_{\mu_k g}(\bz^k) -  \bx^{k+1} = \beta^{-1}(\bz^{k} - \bz^{k+1}),
\end{align}
we obtain
\begin{align}
\calL_{\rho_k,\mu_k}(\bx^{k+1},\bz^{k+1})
 - \calL_{\rho_k,\mu_k}(\bx^{k+1},\bz^{k})
\le -\frac{1}{\mu_k}(\beta^{-1}-\frac{1}{2})
   \|\bz^{k+1}-\bz^{k}\|^{2}. \label{eq:z-update-bound}
\end{align}

For the second part of \eqref{eq:objective-function-decomposition}, by the optimality condition of the subproblem \eqref{eq:bx_update_Polyak} or \eqref{eq:bx_update_storm}, together with its strong convexity, one has
\begin{align}
h(\bx^{k+1})
+ \langle \be^{k} + \nabla Q_{\rho_k}(\bx^{k}),\, \bx^{k+1}-\bx^{k}\rangle
   + \frac{1}{2\mu_k}\|\bx^{k+1}-\bz^{k}\|^{2} \le h(\bx^{k}) + \frac{1}{2\mu_k}\|\bx^{k}-\bz^{k}\|^{2}
   - \frac{1}{2\mu_k}\|\bx^{k+1}-\bx^{k}\|^{2}. \notag
\end{align}
Using the $L_{\rho_k}$-smoothness of $Q_{\rho_k}$ and the fact
$\langle \ba,\bb\rangle \le \tfrac{1}{4\mu_k}\|\ba\|^{2}+ \mu_k\|\bb\|^{2}$, one has
\begin{align}
\calL_{\rho_k,\mu_k}(\bx^{k+1},\bz^{k})
 - \calL_{\rho_k,\mu_k}(\bx^{k},\bz^{k})
\le -\frac{(2\mu_k)^{-1}- L_{\rho_k}}{2}\|\bx^{k+1}-\bx^{k}\|^{2}
+ \mu_k\|\be^{k}\|^{2}.
\label{eq:x-update-bound}
\end{align}
Then, combining \eqref{eq:z-update-bound} and \eqref{eq:x-update-bound} in
\eqref{eq:objective-function-decomposition}, we obtain
\begin{align}
\calL_{\rho_k,\mu_k}(\bw^{k+1}) - \calL_{\rho_k,\mu_k}(\bw^{k})
\le  - \min\{\frac{(2\mu_k)^{-1}-L_{\rho_k}}{2},\,
\frac{1}{\mu_k}(\beta^{-1}-\tfrac{1}{2})\}
\|\bw^{k+1}-\bw^{k}\|^{2}  + \mu_k\|\be^{k}\|^{2}.
\label{eq:combined-update-bound}
\end{align}
Since \eqref{Eq: basic parameter conditions for both algorithms} ensures $\frac{(2\mu_k)^{-1} - L_{\rho_k}}{2} \le \frac{1}{\mu_k}(\beta^{-1}-\tfrac{1}{2})$, one has \( \min\{\frac{(2\mu_k)^{-1}-L_{\rho_k}}{2},\,
\frac{1}{\mu_k}(\beta^{-1}-\tfrac{1}{2})\} = \frac{(2\mu_k)^{-1}-L_{\rho_k}}{2}\). Then, substituting \eqref{eq:penalty-term-bound} and \eqref{eq:combined-update-bound} into
\eqref{eq:function-value-decomposition} yields \eqref{eq:approximate-sufficient-descent}. Since the coefficient of $\| \bw^{k+1}- \bw^{k}\|^{2}$ in \eqref{eq:approximate-sufficient-descent} is positive, rearranging this inequality gives \eqref{eq:sequence-residual-bound}, completing the proof.
\end{proof}

\end{document}
